\chapter{Examen des textes décisifs}\label{ch:textes_decisifs}

Les comptages des chapitres précédents reposent sur un petit nombre de textes, dont la lecture commande les conclusions. Ce chapitre les reprend un à un, sous une forme resserrée : pour chacun, on indique ce qui est lu, ce qui est restitué, ce que dit l’autre témoin, et ce que le texte permet ou ne permet pas de conclure. Le texte complet des deux témoins, la traduction et le commentaire figurent dans les dossiers commentés (chapitre~\ref{ch:dossiers}), auxquels renvoie la page indiquée.

Aucun de ces textes n’a été collationné sur la pierre ou sur une photographie, à l’exception partielle du \nnum{38}, dont un estampage a pu être examiné. Les observations portent sur les lectures transmises par l’EDCS et par l’EDB.

\section{Le nom du Christ et le monogramme}

\begin{description}[style=nextline,leftmargin=1.5em,font=\normalfont\bfseries]
\item[\nnum{38} (352) — Privata, \emph{in pace} + monogramme (p.~\pageref{dos:38})] \emph{Lu} : \emph{bene merenti in pace} ; \emph{EDCS} : \emph{Chr(isti)} ; \emph{EDB} : \emph{((monogramma Christi))} ; estampage : christogramme après IN PACE. \emph{Conclusion} : la pierre porte un monogramme ; sa fonction (abréviation ou signe juxtaposé) reste indéterminée. Le texte n’est pas compté parmi les mentions écrites du Christ.
\item[\nnums{140} (370) et 151 (371) — \emph{in pace Chr(isti)} (p.~\pageref{dos:140}, \pageref{dos:151})] \emph{Lu} : la formule, sous forme abrégée dans les deux bases ; \emph{EDB} : chrisme décrit sur le support. Au \nnum{140}, la formule suit \emph{bene merenti} et précède une seconde mention, \emph{deposita in pace} ; au \nnum{151}, elle clôt le texte après la date. \emph{Conclusion} : formule de repos qui nomme le Christ, probablement au moyen d’un monogramme ; aucun attribut n’est exprimé.
\item[\nnum{311} (388) — \emph{deposita in Chr(isto)} (p.~\pageref{dos:311})] Même situation : abréviation dans les deux bases, chrisme décrit par l’EDB. Le texte est tracé dans le mortier frais. \emph{Conclusion} : comme ci-dessus.
\item[\nnum{458} (401) — \emph{ancilla dei et Chr(ist)i} (p.~\pageref{dos:458})] \emph{Lu} : toute la formule ; l’EDB transcrit \emph{Chr(ist)i}, avec un I final gravé. \emph{Conclusion} : seule abréviation du nom du Christ en lettres assurée ; qualification d’une jeune fille de douze ans, \emph{virgo}, comme servante de Dieu et du Christ. Aucune doctrine des relations divines n’est exprimée.
\end{description}

\section{Les poèmes}

\begin{description}[style=nextline,leftmargin=1.5em,font=\normalfont\bfseries]
\item[\nnum{73} (363) — Eleuteria (p.~\pageref{dos:73})] \emph{Lu dans les deux bases} : \emph{spiritus a carne recedens}, \emph{sanctis pro meritis et opera tanta}, \emph{semper quiescis secura}, \emph{cum lumine claro}, \emph{quae te semper deo dica⟨vi⟩sti}, \emph{ad alta vocaris}. \emph{Restitué et divergent} : le début des vers 1 à 5 (\emph{[at tuus in caelo]} / \emph{[corpore namque tu]us} ; \emph{[vivit cum]} / \emph{[est sociatu]s} ; \emph{[quaeque deu]m metuisti} / \emph{[quaeque de]metuisti}). \emph{Conclusion} : l’esprit de la défunte, séparé de la chair, est associé aux saints en raison de ses mérites, et elle est appelée « vers les hauteurs » ; elle s’était consacrée à Dieu. Le séjour céleste (\emph{in caelo}) et la crainte de Dieu ne sont que des restitutions de l’EDCS.
\item[\nnum{244} (383) — Proiecta, par Damase (p.~\pageref{dos:244})] \emph{Lu} : le poème entier, signé \emph{haec Damasus praestat cunctis solacia fletus}. \emph{Conclusion} : Proiecta, morte à seize ans peu après son mariage, « désirant monter vers la lumière éthérée du ciel » ; le langage est celui de la consolation poétique, sans mention du Christ.
\item[\nnum{306} (389) — Maximilla et Lucceia (p.~\pageref{dos:306})] \emph{Lu dans les deux bases} : \emph{vitam Cristi qui sperat aeternam}, \emph{tunc poterit mors ipsa mori}, \emph{vita perennis erit}, \emph{religiosa vita}, \emph{mirae fidei}. \emph{Restitué} : \emph{[servare foedera C]hristi}. \emph{Divergence} : la ligne \emph{et innocentiae optimique propositi Maximilla virgo ancilla dei civis Pannonia} manque dans le texte de l’EDB. \emph{Conclusion} : le texte le plus explicite du corpus sur la vie éternelle ; la qualification de Maximilla comme \emph{virgo ancilla dei} doit être contrôlée sur l’ICVR V 13355.
\item[\nnum{352} (393) — Eutychius (p.~\pageref{dos:352})] \emph{Lu} : \emph{[- -]ristum credens praemia lucis habet}. \emph{Restitué} : \emph{[qui Ch]-}. \emph{Conclusion} : restitution très probable, contrainte par la finale et le sens ; la foi au Christ et sa récompense (la lumière) sont exprimées dans un poème recueilli dans les \emph{Carmina Latina Epigraphica} (CLE 1346).
\item[\nnum{382} (395) — Panthagate (p.~\pageref{dos:382})] \emph{Lu} : \emph{celestia regna t[enet]} (en partie), \emph{tante fides}, \emph{aurea sunt caelestis quae tu caelo}, \emph{trium subolum mater}. \emph{Restitué} : \emph{in sin[u dei pausat]}. \emph{Conclusion} : le ciel et la fidélité sont lus ; Dieu n’est nommé que dans la restitution.
\end{description}

\section{Dieu, l’Esprit, la Trinité}

\begin{description}[style=nextline,leftmargin=1.5em,font=\normalfont\bfseries]
\item[\nnum{214} (380) — \emph{ancilla dei}, prière pour l’enfant (p.~\pageref{dos:214})] \emph{Lu dans les deux bases} : \emph{hic quiescit ancilla dei}, \emph{pro hunc unum ora subolem quem superstitem reliquisti}. \emph{EDB seule} : trois lignes supplémentaires, dont la date (\emph{Gratiano V et Teodosio}) et deux autres défunts. \emph{Conclusion} : le seul témoignage assuré d’une demande d’intercession adressée à un défunt.
\item[\nnum{249} (383) — l’Esprit reçu (p.~\pageref{dos:249})] \emph{Lu} : \emph{-tum spiritum}, \emph{Mai(as) Merobaud-}. \emph{Restitué} : \emph{ac[cepti / cepit sanc]tum}, le quantième, l’itération du consulat. \emph{Conclusion} : mention probable de l’Esprit saint reçu, peut-être au cours de l’initiation ; le rite et sa date ne peuvent être précisés.
\item[\nnum{400} (397) — Severianus (p.~\pageref{dos:400})] \emph{Lu} : \emph{hic dormit … in somnio pacis}, \emph{cuius spiritus in luce domini susceptus est}. \emph{EDB} : monogrammes et croix monogrammatiques dans le texte. \emph{Conclusion} : l’esprit du défunt est accueilli dans la lumière du Seigneur ; le repos est décrit comme un sommeil de paix.
\item[\nnum{474} (403) — Quintilianus (p.~\pageref{dos:474})] \emph{Lu dans les deux bases} : \emph{homo dei}, \emph{confirmans trinitatem}, \emph{amans castitatem}, \emph{respuens mundum}. \emph{Conclusion} : éloge ascétique qui mentionne la Trinité comme objet de la foi du défunt, sans nommer les personnes.
\item[\nnum{475} (403) — \emph{dedit deus deus tulit} (p.~\pageref{dos:475})] \emph{Lu} : la phrase. \emph{Conclusion} : citation libre de Jb 1, 21, en consolation.
\end{description}

\section{Les textes dont la portée religieuse dépend de restitutions}

\begin{description}[style=nextline,leftmargin=1.5em,font=\normalfont\bfseries]
\item[\nnum{43} (355) — \emph{In nomine [Christi]} (p.~\pageref{dos:43})] Le nom du Christ est entièrement restitué dans les deux bases ; l’EDB décrit un chrisme avec alpha et oméga. L’invocation initiale est vraisemblable ; elle n’est pas lue.
\item[\nnum{232} (382) — \emph{in dom(ino?) cultrix}, \emph{limina sanctorum} (p.~\pageref{dos:232})] \emph{Dom(ino)} est développé avec un point d’interrogation ; \emph{limina sanctorum} est lu, le premier mot de l’expression restitué. Le texte atteste la valeur attachée à une sépulture près des saints ; la désignation de Dieu est douteuse.
\item[\nnum{447} (401) — imprécation grecque (p.~\pageref{dos:447})] Sont lus \emph{ὁρκίζομε δὲ ὑμᾶς}, \emph{-ντοκράτορα}, \emph{εἴ τις το-}, \emph{τὰ ὀστέα αὐτοῦ}, \emph{ἔνοχος τῇ}, \emph{ἡμαίραν τῆς}. Le nom de Dieu est restitué ; la sanction l’est différemment par les deux bases (résurrection, avec un point d’interrogation, dans l’EDCS ; jugement dans l’EDB). Le texte atteste une adjuration par le Tout-Puissant et une menace eschatologique ; il n’atteste ni la résurrection ni le jugement par leur nom.
\end{description}

\section{Bilan}

Les conclusions religieuses du volume reposent sur une quinzaine de textes. Huit d’entre eux sont lus sans restitution dans la partie qui porte l’énoncé religieux (\nnums{214}, 244, 306 pour l’essentiel, 400, 458, 474, 475 et, pour \emph{deo dicavisti}, 73). Les autres dépendent, à des degrés divers, de restitutions ou de développements dont la confrontation des deux bases montre qu’ils ne sont pas uniques. C’est sur ces textes que devrait porter en priorité un contrôle photographique.
